% Vietnamese translation for OI Public Library
% Source-PDF: 英文资料 ENGLISH MATERIALS/Time Bounds For Selection - Blum_Floyd_Pratt_Rivest_Tarjan.pdf
\documentclass[11pt]{article}
\usepackage{oipl-vn}

\newcommand{\PICK}{\textsc{PICK}}
\newcommand{\PICKone}{\textsc{PICK1}}
\newcommand{\PICKonea}{\textsc{PICK1a}}
\newcommand{\PICKoneb}{\textsc{PICK1b}}
\newcommand{\PICKtwo}{\textsc{PICK2}}
\newcommand{\defeq}{\mathrel{\vcentcolon=}}

\title{Các Cận Thời Gian Cho Bài Toán Chọn}
\author{Manuel Blum, Robert W. Floyd, Vaughan Pratt, Ronald L. Rivest, Robert E. Tarjan}
\date{}

\begin{document}
\maketitle

\origtitle{Time Bounds for Selection}
\sourcepath{English Materials/Time Bounds For Selection - Blum\_Floyd\_Pratt\_Rivest\_Tarjan.pdf}

\begin{abstract}
Số phép so sánh cần để chọn phần tử nhỏ thứ \(i\) trong \(n\) số được chứng
minh là bị chặn trên bởi một hàm tuyến tính theo \(n\), thông qua phân tích
một thuật toán chọn mới, \PICK. Cụ thể, không bao giờ cần quá \(5.4305n\) phép
so sánh. Cận này được cải thiện khi \(i\) nằm gần hai đầu của dãy thứ hạng, và
một cận dưới mới cho số phép so sánh cần thiết cũng được chứng minh.
\end{abstract}

\section{Giới thiệu}

Bài báo này trình bày một thuật toán chọn mới, \PICK, và từ phân tích hiệu
quả của nó suy ra kết quả đáng chú ý rằng chi phí chọn bị chặn trên bởi một
hàm tuyến tính theo số phần tử đầu vào. Ngoài ra, bài báo cũng chứng minh một
cận dưới mới cho chi phí chọn.

Bài toán chọn có thể được minh họa rõ nhất bằng việc tính trung vị. Tổng quát
hơn, ta muốn chọn phần tử nhỏ thứ \(i\) trong một tập gồm \(n\) số phân biệt,
hoặc phần tử có thứ hạng gần nhất với một mức phần trăm cho trước.

Sự quan tâm đến bài toán này có thể truy về các cuộc thi đấu thể thao, đặc
biệt là thiết kế giải đấu quần vợt để chọn người giỏi nhất và người giỏi nhì.
Năm 1883, Lewis Carroll công bố một bài viết phê phán cách xác định người
giỏi nhì trong một giải ``loại trực tiếp'': người thua trận chung kết thường
không nhất thiết là người giỏi nhì, vì bất kỳ người chơi nào chỉ thua người
giỏi nhất cũng có thể là người giỏi nhì. Khoảng năm 1930, Hugo Steinhaus đưa
vấn đề này vào phạm vi độ phức tạp thuật toán bằng câu hỏi: cần ít nhất bao
nhiêu trận để xác định đúng cả người giỏi nhất lẫn người giỏi nhì trong \(n\)
đấu thủ? Năm 1932, J. Schreier chứng minh rằng không cần quá
\[
  n+\lceil\log_2 n\rceil-2
\]
trận, và năm 1964, S. S. Kislitsin chứng minh số này cũng là cần thiết.
Phương pháp của Schreier dùng một giải loại trực tiếp để tìm người thắng, rồi
một giải loại trực tiếp thứ hai giữa các đấu thủ đã thua người thắng để tìm
người về nhì.

Với \(i>2\), số trận tối thiểu cần để chọn người giỏi thứ \(i\) trong \(n\)
đấu thủ chỉ được biết với các giá trị nhỏ của \(n\). Thủ tục chọn tổng quát
tốt nhất trước đó là của Hadian và Sobel, cần không quá
\[
  n-i+(i-1)\lceil\log_2(n-i+2)\rceil
\]
trận. Họ tạo một giải loại trực tiếp gồm \(n-i+2\) người, rồi lần lượt loại
\(i-1\) người ``quá giỏi'' để có thể là người giỏi thứ \(i\), bằng kỹ thuật
chọn có thay thế.

Trong tài liệu không có ký hiệu thống nhất cho ``phần tử giỏi thứ \(i\)''.
Bài báo dùng hai toán tử sau.

\begin{itemize}
  \item \(i\circ S\), đọc là ``phần tử thứ \(i\) của \(S\)'', là phần tử nhỏ
  thứ \(i\) của \(S\), với \(1\le i\le |S|\). Độ lớn của \(i\circ S\) tăng
  khi \(i\) tăng. Khi \(S\) đã rõ, ta viết ngắn là \(i\circ\).
  \item \(x\rho S\), đọc là ``hạng của \(x\) trong \(S\)'', là hạng của
  \(x\) trong \(S\), sao cho
  \[
    (x\rho S)\circ S=x.
  \]
\end{itemize}

Gọi \(f(i,n)\) là chi phí tối thiểu trong trường hợp xấu nhất, tức số phép so
sánh nhị phân cần để chọn \(i\circ\), với \(|S|=n\). Ta cũng dùng ký hiệu
\[
  F(\alpha)\defeq
  \limsup_{n\to\infty}
  \frac{f(\lfloor \alpha(n-1)\rfloor+1,n)}{n},
  \qquad 0\le\alpha\le 1,
\]
để đo độ khó tương đối của việc tính các mức phần trăm.

Trong Phần 2, kết quả chính \(f(i,n)=O(n)\) được chứng minh bằng cách phân tích
thuật toán chọn cơ bản \PICK. Trong Phần 3, \PICK{} được tinh chỉnh để thu được
các cận chặt nhất của bài báo:
\[
  \max_{0\le\alpha\le 1} F(\alpha)\le 5.4303
  \tag{1}
\]
và
\[
  F(\alpha)\le
  1+4.4303\frac{\alpha}{\beta}
  +10.861\alpha\lceil\log_2(\beta/\alpha)\rceil,
  \qquad 0<\alpha<\beta,
  \tag{2}
\]
trong đó \(\beta=0.203688\ldots\). Trong Phần 4, bài báo suy ra cận dưới
\[
  F(\alpha)\ge 1+\min(\alpha,1-\alpha),
  \qquad 0\le\alpha\le 1.
  \tag{3}
\]
Không có bằng chứng nào cho thấy các bất đẳng thức (1)--(3) là tốt nhất có thể;
tác giả phỏng đoán rằng chúng còn có thể được cải thiện đáng kể.

\section{Thuật toán chọn mới \PICK}

Trong phần này ta mô tả thuật toán cơ bản và chứng minh \(f(i,n)=O(n)\). Giả
sử cần chọn \(i\circ S\), với \(|S|=n\).

\PICK{} hoạt động bằng cách lần lượt loại khỏi \(S\) các tập con mà mọi phần tử
trong đó đã biết là quá lớn hoặc quá nhỏ để có thể là \(i\circ\), cho đến khi
chỉ còn lại \(i\circ\). Mỗi tập con bị loại chứa ít nhất một phần tư số phần
tử còn lại. \PICK{} khá giống thuật toán \textsc{Find} của Hoare, ngoại trừ
việc phần tử \(m\) dùng để phân hoạch \(S\) được chọn cẩn thận hơn.

Thuật toán được mô tả bằng ba hàm phụ \(b(i,n)\), \(c(i,n)\), và \(d(i,n)\);
khi không gây nhầm lẫn, ta bỏ qua danh sách đối số của chúng. Vì trọng tâm là
tính chất tiệm cận của \PICK, các chi tiết cho trường hợp \(n\bmod c\ne 0\) sẽ
được bỏ qua ở đây và được xử lý trong Phần 3.

\paragraph{Thuật toán \PICK.}
Chọn \(i\circ S\), với \(|S|=n\) và \(1\le i\le n\).

\begin{enumerate}
  \item Chọn một phần tử \(m\in S\).
  \begin{enumerate}
    \item Sắp \(S\) thành \(n/c\) cột, mỗi cột có độ dài \(c\), rồi sắp xếp
    từng cột.
    \item Gọi \(T\) là tập gồm \(n/c\) phần tử, mỗi phần tử là phần tử nhỏ
    thứ \(d\) của một cột. Chọn
    \[
      m=b\circ T.
    \]
    Nếu \(n/c>1\), dùng đệ quy \PICK{} để chọn phần tử này.
  \end{enumerate}
  \item Tính \(m\rho S\): so sánh \(m\) với mọi phần tử \(x\in S\) mà ta chưa
  biết \(m<x\) hay \(m>x\).
  \item Nếu \(m\rho S=i\), dừng, vì \(m=i\circ S\). Nếu \(m\rho S>i\), loại
  \[
    D=\{x\mid x\ge m\}
  \]
  và đặt \(n\gets n-|D|\). Nếu \(m\rho S<i\), loại
  \[
    D=\{x\mid x\le m\},
  \]
  rồi đặt \(n\gets n-|D|\) và \(i\gets i-|D|\). Quay lại bước 1.
\end{enumerate}

Sau bước 1, thứ tự bộ phận đã xác định trên \(S\) có thể hình dung như sau.
Các cột có phần tử lớn nhất ở trên cùng. Tập \(T\) gồm các phần tử nhỏ thứ
\(d\) của các cột. Vì lời gọi đệ quy ở bước 1(b) đã xác định phần tử nào của
\(T\) nhỏ hơn \(m\), phần tử nào lớn hơn \(m\), các cột được chia thành hai
nhóm: \(b-1\) cột có phần tử nhỏ thứ \(d\) nhỏ hơn \(m\), và \(n/c-b\) cột có
phần tử nhỏ thứ \(d\) lớn hơn \(m\). Mọi phần tử trong vùng \(G\) chắc chắn lớn
hơn \(m\), còn mọi phần tử trong vùng \(L\) chắc chắn nhỏ hơn \(m\). Vì vậy, ở
bước 2 chỉ cần so sánh \(m\) với các phần tử trong hai vùng còn chưa quyết
định, gọi là \(A\) và \(B\).

Không có phần tử nào bị loại sai ở bước 3. Nếu \(m\rho S>i\) thì \(m\) quá
lớn, nên \(m\) và mọi phần tử lớn hơn nó đều có thể bị loại; trường hợp
\(m\rho S<i\) đối xứng. Chú ý rằng ít nhất toàn bộ \(G\) hoặc toàn bộ \(L\) sẽ
bị loại.

\paragraph{Định lý 1.}
\[
  f(i,n)=O(n).
\]

\paragraph{Chứng minh.}
Ta chỉ ra rằng một lựa chọn hợp lý cho \(b(i,n)\), \(c(i,n)\), và \(d(i,n)\)
cho một thuật toán chọn thời gian tuyến tính. Gọi \(h(c)\) là chi phí sắp xếp
\(c\) số bằng thuật toán Ford--Johnson. Đã biết rằng
\[
  h(c)=\sum_{1\le j\le c}\left\lceil\log_2\frac{3j}{4}\right\rceil.
  \tag{4}
\]
Chi phí của bước 1(a) là \(n\cdot h(c)/c\), từ đó thấy ngay rằng \(c(i,n)\)
phải bị chặn trên bởi một hằng số nếu muốn \PICK{} chạy trong thời gian tuyến
tính.

Gọi \(P(n)\) là chi phí lớn nhất của \PICK{} trên mọi \(i\). Chi phí của bước
1(b) bị chặn bởi \(P(n/c)\). Từ cấu trúc các vùng \(L\) và \(G\) ở trên, ta có
\[
  P(n)\le
  n\frac{h(c)}{c}
  +P\!\left(\frac{n}{c}\right)
  +n
  +P(n-\min(|L|,|G|)).
  \tag{5}
\]

Để tối ưu hóa chặn này, chọn \(c=21\), \(d=11\), và \(b=n/(2c)=n/42\). Khi đó
\(m\) là trung vị của \(T\), còn \(T\) là tập các trung vị cột. Vì
\(h(21)=66\), suy ra
\[
  P(n)\le \frac{66n}{21}
  +P\!\left(\frac{n}{21}\right)
  +n
  +P\!\left(\frac{31n}{42}\right).
  \tag{6}
\]
Bằng quy nạp toán học, điều này kéo theo
\[
  P(n)\le \frac{58n}{3}<19.4n.
  \tag{7}
\]
Trường hợp cơ sở dựa trên việc \(h(n)<19n\) với \(n<10^5\), nên các trường
hợp nhỏ có thể xử lý bằng cách sắp xếp. \PICK{} chạy thời gian tuyến tính vì
mỗi lượt loại được một phần đáng kể của \(S\), với chi phí tỉ lệ với số phần
tử còn lại trong lượt đó. Chú ý rằng cần \(c\ge 5\) để \PICK{} chạy trong
thời gian tuyến tính.

\section{Cải thiện \PICK}

Kết quả chính \(f(i,n)=O(n)\) đã được chứng minh. Phần này phân tích chi tiết
các phiên bản cải tiến của \PICK. Bài báo mô tả hai sửa đổi: \PICKone, cho cận
tổng quát tốt nhất đối với \(F(\alpha)\), và \PICKtwo, hiệu quả hơn \PICKone{}
khi \(i<\beta n\) hoặc \(i>(1-\beta)n\), với \(\beta=0.203688\ldots\). Nhờ
tính đối xứng, trong suốt phần này giả sử \(i\le \lfloor n/2\rfloor\).

\paragraph{Định lý 2.}
Với \(0\le\alpha\le 1\),
\[
  F(\alpha)\le 5.4303.
\]

\paragraph{Chứng minh.}
Ta phân tích \PICKone, khác \PICK{} ở các điểm sau.

\begin{enumerate}
  \item Các phần tử của \(S\) chỉ được sắp vào cột một lần. Sau đó, những cột
  bị phá vỡ bởi thao tác loại bỏ được khôi phục về độ dài đầy đủ bằng một
  bước trộn mới ở cuối mỗi lượt.
  \item Bước phân hoạch được sửa để số phép so sánh là một hàm tuyến tính của
  số phần tử cuối cùng bị loại.
  \item Thao tác loại bỏ không phá vỡ quá nửa số cột trong mỗi lượt, nhờ đó
  các sửa đổi khác hoạt động tốt.
  \item Bước sắp xếp ngầm trong lời gọi đệ quy để chọn \(m\) được thay một
  phần bằng bước trộn kể từ lượt thứ hai trở đi, vì điểm trước bảo đảm rằng
  một nửa tập \(T\) ở lượt \(j\) cũng đã xuất hiện trong lời gọi đệ quy ở lượt
  \(j-1\).
\end{enumerate}

Một \term{\(k\)-cột} là một cột đã sắp có độ dài \(k\). Giá trị tối ưu của
hàm \(c\) là \(15\), và sẽ được dùng trực tiếp.

\paragraph{Thủ tục \PICKone.}
Chọn \(i\circ S\), với \(|S|=n\) và \(1\le i\le\lfloor n/2\rfloor\).

\begin{enumerate}
  \item Nếu \(n\le 45\), sắp \(S\), in \(i\circ S\), rồi dừng.
  \item Sắp \(S\) thành \(\lfloor n/15\rfloor\) cột \(15\)-phần tử, và có thể
  thêm một cột gồm \(n\bmod 15\) phần tử.
  \item Dùng thủ tục \PICKonea{} để chọn \(i\circ\).
\end{enumerate}

\paragraph{Thủ tục \PICKonea.}
Giống \PICKone, ngoại trừ việc \(S\) đã được sắp thành các cột \(15\)-phần tử.

\begin{enumerate}
  \item Nếu \(n\le 45\), sắp \(S\), in \(i\circ S\), rồi dừng.
  \item Sắp tập \(T\) gồm các trung vị cột thành các cột \(15\)-phần tử, và có
  thể thêm một cột gồm \(\lfloor n/15\rfloor\bmod 15\) phần tử.
  \item Dùng thủ tục \PICKoneb{} để chọn \(i\circ S\).
\end{enumerate}

\paragraph{Thủ tục \PICKoneb.}
Giống \PICKonea, ngoại trừ việc \(T\) cũng đã được sắp thành các cột
\(15\)-phần tử.

\begin{enumerate}
  \item Dùng \PICKonea{} để chọn \(m\), trung vị của \(T\).
  \item Phân hoạch \(A\cup B\) quanh \(m\), và dừng ngay khi đã rõ
  \(m\rho S<i\) hoặc \(m\rho S>i\):
  \begin{enumerate}
    \item Chèn \(m\) vào mỗi \(7\)-cột của \(B\) bằng chèn nhị phân, tốn \(3\)
    phép so sánh mỗi cột.
    \item Chèn \(m\) vào mỗi \(7\)-cột của \(A\) bằng tìm kiếm tuyến tính bắt
    đầu gần trung vị của từng cột \(15\)-phần tử.
  \end{enumerate}
  \item Nếu \(m\rho S=i\), in \(m=i\circ S\), rồi dừng. Nếu \(m\rho S>i\), loại
  \[
    G\cup\{x\mid x\in B,\ x>m\}.
  \]
  Nếu \(m\rho S<i\), loại
  \[
    L\cup\{x\mid x\in A,\ x<m\},
  \]
  rồi giảm \(i\) đi số phần tử bị loại.
  \item Khôi phục \(S\) thành một tập các cột \(15\)-phần tử bằng các phép
  trộn sau. Gọi \(|X|\) là số phần tử của một tập \(X\). Gọi \(U\) là tập các
  cột có độ dài nhỏ hơn \(15\); thực tế, mỗi cột của \(U\) có độ dài không quá
  \(7\). Gọi \(Y\subset U\) là tập các cột ngắn nhất của \(U\), sao cho
  \[
    |Y|=|U|\bmod 15,
  \]
  và gọi \(V\) là tập mọi \(7\)-cột trong \(U-Y\). Chia
  \(U-(V\cup Y)\) thành hai tập con \(X\) và \(W\), sao cho \(W\) chứa \(w\)
  cột, các cột của \(W\) không ngắn hơn các cột của \(X\), và
  \[
    |W|+|X|=7w.
  \]
  Sau đó:
  \begin{enumerate}
    \item Kéo dài mọi cột trong \(W\) lên độ dài \(7\) bằng cách dùng chèn nhị
    phân để đặt từng phần tử của \(X\) vào một cột của \(W\).
    \item Khi đó mọi cột trong \(U-Y\) đều là \(7\)-cột. Trộn chúng từng cặp
    để tạo thành các \(14\)-cột.
    \item Dùng chèn nhị phân để đặt từng phần tử của \(Y\) vào một \(14\)-cột.
    Khi đó \(S\) đã được khôi phục thành một tập các cột \(15\)-phần tử.
  \end{enumerate}
  \item Khôi phục tập \(T\) các trung vị cột thành các cột \(15\)-phần tử. Gọi
  \(Z\subset T\) là các trung vị cột đã là trung vị cột ở bước 1. Các phần tử
  của \(Z\) đã được sắp thành các cột có kích thước ít nhất \(8\), vì bước 3
  của lời gọi đệ quy ở bước 1 đã loại \(Z\) theo các chuỗi có kích thước như
  vậy.
  \begin{enumerate}
    \item Trộn các cột của \(Z\) để tạo thành các cột \(15\)-phần tử và một số
    phần dư, xử lý riêng từng kích thước cột:
    \[
    \begin{array}{rcl}
    2(8)  &:& 8+7,\quad 1\text{ phần dư},\\
    5(9)  &:& 9+6,\quad 9+6,\quad 9+3+3,\quad 0\text{ phần dư},\\
    3(10) &:& 10+5,\quad 10+5,\quad 0\text{ phần dư},\\
    4(11)+1(1) &:& 11+4,\quad 11+4,\quad 11+3+1,\quad
                 0\text{ phần dư}.
    \end{array}
    \]
    Với \(11\)-cột, trước tiên tách riêng \(1/45\) số \(11\)-cột và phá chúng
    thành các \(1\)-cột. Với các cột độ dài \(12\) trở lên, tách riêng một số
    phần tử để chèn nhị phân vào các cột còn lại. Cuối cùng, sắp các phần dư
    thành các cột \(15\)-phần tử.
    \item Sắp \(T-Z\) thành các cột \(15\)-phần tử. Lúc này \(T\) đã được khôi
    phục thành một tập các cột \(15\)-phần tử.
  \end{enumerate}
  \item Giảm \(n\) đi số phần tử bị loại ở bước 5. Nếu \(n\le 45\), sắp \(S\),
  in \(i\circ S\), rồi dừng; nếu không, quay lại bước 1.
\end{enumerate}

Để phân tích \PICKone, dùng các ký hiệu sau:
\[
\begin{array}{rcl}
P_1(n),P_{1a}(n),P_{1b}(n)
&=&\text{chi phí lớn nhất của }\PICKone,\PICKonea,\PICKoneb,\\
v&=&\text{số phép so sánh trong bước }\PICKoneb(2\mathrm{ii}),\\
d&=&\text{số phần tử từ }A\cup B\text{ bị loại ở bước }\PICKoneb(3),\\
g_a,g_b&=&\text{số phần tử từ }A,B\text{ được thấy là lớn hơn }m,\\
l_a,l_b&=&\text{số phần tử từ }A,B\text{ được thấy là nhỏ hơn }m,\\
w,x,y&=&\text{số cột trong các tập }W,X,Y\text{ ở bước }\PICKoneb(4).
\end{array}
\]
Vì \(h(15)=42\), ta có ngay
\[
  P_1(n)\le \frac{42n}{15}+P_{1a}(n)
  =2.8n+P_{1a}(n),
  \tag{8}
\]
và
\[
  P_{1a}(n)\le \frac{42n}{225}+P_{1b}(n)
  =0.186\overline{6}n+P_{1b}(n).
  \tag{9}
\]

\paragraph{Bổ đề 1.}
\[
  v\le d+\frac{n}{30}.
  \tag{10}
\]

\paragraph{Chứng minh.}
Có hai trường hợp, tùy theo bước \PICKoneb(3) loại \(L\) hay \(G\).

Nếu \(L\) bị loại, thì rõ ràng có nhiều nhất một phép so sánh cho mỗi cột
trong \(A\), cộng thêm một phép cho mỗi phần tử của \(A\) bị loại.

Nếu \(G\) bị loại, thì từ điều kiện dừng của phân hoạch và do \(i\le
\lfloor n/2\rfloor\), suy ra số phần tử của \(B\) được chứng minh lớn hơn
\(m\) không nhỏ hơn số phần tử tương ứng cần để trả cho các phép so sánh tuyến
tính trong \(A\). Nói cách khác, \(g_b\ge l_a\), còn \(l_a\ge v-n/30\) như
trường hợp trước. Vì trong trường hợp này \(d=g_b\), ta được bổ đề.

\paragraph{Bổ đề 2.}
\[
  |X|\le \frac{6d}{7}.
  \tag{11}
\]

\paragraph{Chứng minh.}
Từ cách chọn \(X,W,Y\),
\[
  d\ge 7(w+x+y)-|W|-|X|-|Y|,
\]
và \(7w-|W|=|X|\). Do đó
\[
  d\ge 7x+7y-|Y|\ge 7x.
\]
Vì \(6x\ge |X|\), suy ra \(d\ge 7|X|/6\).

Bước \PICKoneb{}(5i) tốn trong trường hợp xấu nhất \(21/20\) phép so
sánh trên mỗi phần tử để trộn và sắp \(Z\) thành các cột \(15\)-phần tử; trường
hợp xấu này xảy ra khi \(Z\) chỉ chứa các \(8\)-cột. Vì \(|Z|=n/30\), bước này
tốn nhiều nhất \(7n/200\) phép so sánh.

Từ đó có truy hồi
\[
\begin{aligned}
P_{1b}(n)\le{}&
P_{1a}(\lfloor n/15\rfloor)
 +3(n/30)
 +(d+n/30)
 +3(6d/7)\\
&+\frac{13(7n/30-d)}{15}
 +\frac{4(7n/30-d)}{15}
 +\frac{7n}{200}\\
&+\frac{42(7n/30-d)}{225}
 +P_{1b}(11n/15-d).
\end{aligned}
\]
Các hạng tử lần lượt ứng với các bước 1, \(2\mathrm{i}\), \(2\mathrm{ii}\),
\(4\mathrm{i}\), \(4\mathrm{ii}\), \(4\mathrm{iii}\), \(5\mathrm{i}\),
\(5\mathrm{ii}\), và các lượt tiếp theo. Rút gọn cho
\[
  P_{1b}(n)\le (n+5d)
  +P_{1b}\!\left(\frac{13197n}{27000}-d\right).
  \tag{12}
\]
Vế phải của (12) lớn nhất tại \(d=0\), nên
\[
  P_{1b}(n)\le \frac{13197n}{27000-13197}<2.443n,
  \tag{13}
\]
\[
  P_{1a}(n)<2.6305n,
  \tag{14}
\]
và
\[
  P_1(n)<5.4305n.
  \tag{15}
\]
Vì
\[
  \max_{1\le c\le 45}\frac{h(c)}{c}
  =\frac{h(45)}{45}<5.43,
  \tag{16}
\]
cơ sở quy nạp cho (12) là hợp lệ; điều này cũng xử lý các bước
\PICKone(1), \PICKonea(1), và \PICKoneb(6). Định lý được chứng minh.

\paragraph{Định lý 3.}
Với \(0<\alpha\le\beta\),
\[
  F(\alpha)\le
  1+4.4303\frac{\alpha}{\beta}
  +10.861\alpha\lceil\log_2(\beta/\alpha)\rceil,
  \tag{17}
\]
trong đó \(\beta=0.203688\ldots\).

\paragraph{Chứng minh.}
Ta phân tích \PICKtwo, về bản chất giống \PICK{} với
\[
  b(i,n)=i,\qquad c(i,n)=2,\qquad d(i,n)=1,
\]
và bỏ bước phân hoạch.

\paragraph{Thủ tục \PICKtwo.}
Chọn \(i\circ S\), với \(|S|=n\) và \(i<\beta n\).

\begin{enumerate}
  \item So sánh các phần tử của \(S\) từng cặp, tạo \(\lfloor n/2\rfloor\)
  cặp và có thể một phần tử lẻ.
  \item Gọi \(T\) là tập các phần tử nhỏ hơn trong mỗi cặp. Nếu
  \(i<\beta\lfloor n/2\rfloor\), dùng \PICKtwo; nếu không, dùng \PICKone{} để
  chọn \(m\), phần tử nhỏ thứ \(i\) của \(T\).
  \item Loại mọi phần tử đã biết là lớn hơn \(m\). Cụ thể, trong các cặp mà
  phần tử nhỏ hơn đã lớn hơn \(m\), cả hai phần tử đều lớn hơn \(m\) và bị
  loại.
  \item Dùng \PICKone{} để chọn \(i\circ S\) từ phần còn lại của \(S\).
\end{enumerate}

Khi \(i=1\), thủ tục này trở thành một giải loại trực tiếp đơn giản. Gọi
\(P_2(i,n)\) là số phép so sánh lớn nhất mà \PICKtwo{} dùng để chọn \(i\circ\).
Ta có
\[
  P_2(i,n)\le
  \lfloor n/2\rfloor
  +\min\{P_1(\lfloor n/2\rfloor),P_2(i,\lfloor n/2\rfloor)\}
  +P_1(2i).
  \tag{18}
\]
Với các giá trị cụ thể của \(i\) và \(n\), thủ tục \PICKtwo{} được gọi liên
tiếp
\[
  t=\lceil\log_2(\beta n/i)\rceil
\]
lần trong các lời gọi đệ quy ở bước 2 trước khi \PICKone{} được gọi. Vì vậy
\[
  P_2(i,n)\le
  \sum_{0\le j<t}\left\lfloor\frac{n}{2^{j+1}}\right\rfloor
  +P_1\!\left(\frac{n}{2^t}\right)
  +tP_1(2i).
  \tag{19}
\]
Bất đẳng thức này trực tiếp suy ra định lý. Giá trị đúng của
\(\beta=0.203688\ldots\) là giá trị lớn nhất sao cho
\[
  P_2(\lfloor\beta n\rfloor,n)<P_1(n).
\]

Tóm lại, đường cận trên đối với \(F(\alpha)\) dùng \PICKtwo{} khi
\(\alpha\) gần \(0\) hoặc gần \(1\), và dùng \PICKone{} ở vùng giữa. Bài báo
nhận xét rằng không vô lý khi phỏng đoán đường cong thật của \(F(\alpha)\) là
đơn đỉnh và đạt cực đại tại
\(\alpha=1/2\). Độ phức tạp tương đối của \PICKone{} cũng khiến các tác giả
phỏng đoán rằng cận trên còn có thể được cải thiện đáng kể.

\section{Một cận dưới}

Phần này suy ra một cận dưới cho \(F(\alpha)\) bằng phương pháp ``đối thủ''.
Knuth gọi kỹ thuật này là xây dựng một ``oracle''. Quá trình chọn có thể được
mô hình hóa như một trò chơi giữa thuật toán chọn, gọi là người chơi \(A\), và
đối thủ của nó, gọi là người chơi \(B\). Người chơi \(A\) muốn tìm \(i\circ S\)
với ít phép so sánh nhất; người chơi \(B\) muốn buộc \(A\) dùng nhiều phép so
sánh nhất. Hai người đi luân phiên. Mỗi nước đi của \(A\) là một câu hỏi so
sánh, chẳng hạn ``\(x<y\) hay không?'' với \(x,y\in S\). Ở lượt của mình,
\(B\) phải trả lời ``có'' hoặc ``không''. Các câu trả lời của \(B\) có thể tùy
ý, miễn là không mâu thuẫn với các câu trả lời trước đó. Khi \(A\) đã rút ra
đủ thông tin để xác định \(i\circ S\), trò chơi kết thúc.

Lợi thế của cách nhìn này là một cận dưới không tầm thường cho độ dài trò
chơi, độc lập với chiến lược của \(A\), có thể thu được chỉ bằng cách đưa ra
một chiến lược đủ khéo cho \(B\). Độ dài trò chơi ở đây chính là chi phí
minimax thông thường:
\[
  f(i,n)\defeq \min_A\max_B c(A,B),
  \tag{20}
\]
trong đó \(c(A,B)\) là độ dài trò chơi giữa hai chiến lược cụ thể \(A\) và
\(B\).

Trong trò chơi này, \(B\) đóng vai trò của dữ liệu. Một chiến lược của \(B\)
thực chất là một quy tắc tính ra một bộ dữ liệu đặc biệt xấu cho chiến lược
của \(A\), vì luôn có thể dựng một tập số thật phù hợp với các câu trả lời của
\(B\). Do đó một chiến lược tốt cho \(B\) là một thủ tục làm khó bất kỳ chiến
lược nào của \(A\).

Ta mô tả chiến lược cụ thể của \(B\) để thu được cận dưới. Khi trò chơi diễn
ra, sẽ có nhiều phần tử \(x\) mà \(A\) đã biết đủ để kết luận \(x\ne i\circ\),
tức \(x<i\circ\) hoặc \(x>i\circ\). Ban đầu, \(B\) xem mọi phần tử \(x\in S\)
là thành viên của tập \(U\), nghĩa là \(B\), và do đó cả \(A\), còn chưa chắc
liệu \(x<i\circ\), \(x=i\circ\), hay \(x>i\circ\). Sau một thời gian, \(A\) có
thể buộc \(B\) quyết định trạng thái của một phần tử cụ thể \(x\in U\). Nếu
\(B\) quyết định \(x>i\circ\), \(B\) bỏ \(x\) khỏi \(U\) và đưa vào tập \(G\).
Tương tự, nếu \(B\) quyết định \(x<i\circ\), \(B\) bỏ \(x\) khỏi \(U\) và đưa
vào tập \(L\). Ban đầu \(G\) và \(L\) đều rỗng. Phần tử cuối cùng là
\(i\circ\) sẽ là một phần tử nào đó còn trong \(U\), nên chừng nào \(|U|>1\)
thì trò chơi còn phải tiếp tục. Chiến lược của \(B\) cố giữ \(U\) càng lớn
càng lâu càng tốt.

Theo chiến lược của \(B\), trò chơi gồm hai pha. Chừng nào
\[
  |L|<i-1
  \qquad\text{và}\qquad
  |G|<n-i,
\]
\(B\) còn toàn quyền đưa một phần tử \(x\in U\) vào \(L\) hoặc \(G\). Sau khi
một trong hai tập \(L\) hoặc \(G\) đầy, \(B\) bị ràng buộc hơn nhiều, vì không
được phép làm \(|L|\ge i\) hoặc \(|G|\ge n-i+1\). Khi đó trò chơi suy biến
thành việc \(A\) chỉ còn phải tìm phần tử nhỏ nhất hoặc lớn nhất của \(U\).

Trong pha thứ nhất, \(B\) không bao giờ bỏ quá một phần tử khỏi \(U\) trong
một lượt. Điều này không gây rắc rối nếu phần tử đó là cực đại hoặc cực tiểu
của \(U\), và \(B\) đưa nó vào \(G\) hoặc \(L\) tương ứng. Mỗi phần tử đưa vào
\(G\) được giả sử là nhỏ hơn mọi phần tử đã đưa vào \(G\) trước đó, đồng thời
lớn hơn mọi phần tử còn trong \(U\) và \(L\). Đối xứng, mỗi phần tử đưa vào
\(L\) được giả sử là lớn hơn mọi phần tử đã đưa vào \(L\) trước đó, đồng thời
nhỏ hơn mọi phần tử còn trong \(U\) và \(G\). Quy tắc này xác định hoàn toàn
câu trả lời của \(B\), ngoại trừ khi \(A\) muốn so sánh hai phần tử
\(x,y\in U\). Ngoài ra, \(B\) chỉ bỏ một phần tử khỏi \(U\) khi \(A\) đưa ra
một yêu cầu như vậy.

\(B\) luôn hạn chế \(U\) sao cho mỗi phần tử \(x\in U\) là cực đại hoặc cực
tiểu của \(U\), hoặc cả hai, đối với thứ tự bộ phận đã được các câu trả lời
trước đó của \(B\) cố định trên \(S\). Thậm chí, \(B\) duy trì điều kiện mạnh
hơn: với mỗi \(x\in U\), có nhiều nhất một \(y\in U\) sao cho đã biết
\(x<y\) hoặc \(y<x\). Vì vậy \(U\) chỉ có ba trạng thái phần tử. Đặt
\(\varepsilon(x)\) bằng \(-1\), \(0\), hoặc \(1\), tùy theo \(x\) là phần tử
nhỏ hơn trong một cặp, là phần tử cô lập, hay là phần tử lớn hơn trong một
cặp.

Chiến lược của \(B\) khi so sánh hai phần tử \(x,y\in U\) là như sau; không
mất tính tổng quát, giả sử \(\varepsilon(x)\le\varepsilon(y)\).

\begin{enumerate}
  \item Trả lời rằng \(x\) nhỏ hơn \(y\).
  \item Nếu \(\varepsilon(x)=\varepsilon(y)=0\), không làm gì thêm. Nếu
  \(\varepsilon(x)=-1\), bỏ \(x\) khỏi \(U\) và đưa vào \(L\). Ngược lại, bỏ
  \(y\) khỏi \(U\) và đưa vào \(G\).
\end{enumerate}

Về bản chất, chiến lược của \(B\) tạo một cặp mới trong \(U\) khi
\(\varepsilon(x)=\varepsilon(y)=0\); trong các trường hợp khác, một phần tử bị
bỏ khỏi \(U\), và số cặp trong \(U\) giảm một. Gọi
\[
  c=\text{số phép so sánh đã thực hiện},\qquad
  p=\text{số cặp hiện có trong }U.
\]
Dễ kiểm tra rằng, chừng nào trò chơi còn ở pha thứ nhất, chiến lược của \(B\)
duy trì điều kiện
\[
  c-p+2|U|=2n.
  \tag{21}
\]
Điều này đúng lúc bắt đầu, khi \(c=p=0\) và \(|U|=n\).

Khi pha thứ nhất kết thúc, \(L\) hoặc \(G\) đã đầy, nên
\[
  |U|\ge n-\min(i-1,n-i).
  \tag{22}
\]
Hơn nữa, trong pha thứ hai, \(A\) cần ít nhất \(|U|-1-p\) phép so sánh nữa để
kết thúc trò chơi: tối thiểu \(A\) phải làm công việc tìm phần tử nhỏ nhất
hoặc lớn nhất của \(U\), vốn cần \(|U|-1\) phép so sánh, trong đó \(p\) phép đã
được thực hiện. Vì vậy tổng số phép so sánh ít nhất là
\[
  f(i,n)\ge c+|U|-1-p
  \ge n+\min(i-1,n-i)-1,
  \qquad 1\le i\le n.
  \tag{23}
\]
Cho \(n\to\infty\) với \(i=\lfloor\alpha(n-1)\rfloor+1\), ta được
\[
  F(\alpha)\ge 1+\min(\alpha,1-\alpha).
  \tag{24}
\]
Cận này cũng được vẽ trong hình tổng kết của bản gốc.

\section{Tóm tắt}

Kết quả quan trọng nhất của bài báo là: trong trường hợp xấu nhất, bài toán
chọn có thể thực hiện trong thời gian tuyến tính. Không cần quá \(5.4305n\)
phép so sánh để chọn phần tử nhỏ thứ \(i\) trong \(n\) số, với mọi
\(1\le i\le n\). Cận này có thể cải thiện khi \(i\) gần hai đầu khoảng giá trị
của nó.

Bài báo cũng đưa ra một cận dưới tổng quát, đặc biệt cho thấy cần ít nhất
\[
  \frac{3n}{2}-2
\]
phép so sánh để tính trung vị. Các tác giả tin rằng các hằng số tỉ lệ trong cả
cận trên lẫn cận dưới đều có thể được cải thiện đáng kể.

\section*{Tài liệu tham khảo}

\begin{enumerate}
  \item Lewis Carroll, ``Lawn Tennis Tournaments,'' \textit{St. James's Gazette}
  (August 1, 1883), pp. 5--6. Reprinted in \textit{The Complete Works of Lewis
  Carroll}, New York Modern Library, 1947.
  \item L. R. Ford and S. M. Johnson, A tournament problem, \textit{The American
  Mathematical Monthly} 66 (May 1959), 387--389.
  \item Abdollah Hadian, Optimality properties of various procedures for ranking
  \(n\) different numbers using only binary comparisons, Technical Report 117,
  Dept. of Statistics, Univ. of Minnesota, May 1969, Ph.D. thesis.
  \item Abdollah Hadian and Milton Sobel, Selecting the \(t\)-th largest using
  binary errorless comparisons, Technical Report 121, Dept. of Statistics, Univ.
  of Minnesota, May 1969.
  \item C. A. R. Hoare, Find (Algorithm 65), \textit{Communications of the ACM},
  July 1961, pp. 321--322.
  \item S. S. Kislitsin, On the selection of the \(k\)-th element of an ordered
  set by pairwise comparisons, \textit{Sibirsk. Math. Z.} 5 (1964), 557--564
  (MR 29, No. 2198), Russian.
  \item Donald E. Knuth, \textit{The Art of Computer Programming}, Vol. III:
  Sorting and Searching, Addison-Wesley, 1973.
  \item Józef Schreier, O systemach eliminacjii w turniejach (On elimination
  systems in tournaments), \textit{Mathesis Polska} 7 (1932), 154--160, Polish.
\end{enumerate}

\end{document}
